\section{Discussion}

The central lesson of this project was not scientific but procedural: an appealing
interactive 3D atlas and a FAIR scientific resource are not the same amount of work, and
the difference is not visual polish. rice-atlas already had the harder-to-build part -- a
working, organ-selectable procedural plant viewer. What it lacked was cheaper in
engineering terms but slower in practice: a cited dataset behind every organ, a resolvable
license, and a record of what had and had not been checked. Recreating the visual
experience for \textit{Arabidopsis} took a fraction of the effort spent sourcing,
verifying, and honestly scoping the data behind it.

That verification effort surfaced problems worth naming. A CrossRef check was needed
because a plausible-looking DOI is not evidence it is correct. A live API check against
NASA OSDR was needed because a study accession copied from memory is a guess, not a fact,
until confirmed against the record NASA actually serves. And checking travadb.org's actual
interface -- rather than assuming a REST-style URL pattern would work -- avoided shipping a
broken or misleading link in a resource whose entire premise is that its links and numbers
are trustworthy.

The decision to compute only a simple mean-CPM ratio, rather than a full
differential-expression model, was a scope choice made explicit rather than hidden: a
half-rigorous statistical claim dressed as a rigorous one would be worse than a clearly
labelled simple one. The same logic applied to organ coverage -- three of five organs carry
no spaceflight layer at all, and the viewer says so, rather than reusing the one available
whole-seedling comparison across organs it was never measured in.

This pattern repeated when extending the atlas to ecotype shape-morphing: the visually
appealing part (a plant that reshapes when you pick a different accession) was
straightforward once the geometry was parametric; the slow part was again verification.
The raw-data archive behind Camargo et al.\ \citeyearpar{Camargo2014}'s rosette-shape
descriptors turned out to be genuinely reusable -- a real, column-labelled CSV, not just a
figure to eyeball -- but only after an explanatory table in the same paper's own
supplementary material (numbered, unlabelled columns, six example rows) turned out not to
be it. The lesson generalizes: "supplementary data exists" and "supplementary data is the
specific table you need, in the format you need" are different claims, and only checking
the second one is real verification.

Building the growth animation surfaced the same lesson from a different angle: the two
real bugs that actually blocked it (a crashing Metal shader compiler; a materials helper
silently erasing the ground's texture every frame) were both found only by running the full
pipeline to completion and checking its output, not by reading the code and judging it
correct. A short test render of a handful of frames would not have caught the shader-cache
crash, which took roughly nine frames of repeated recompilation to trigger -- "it worked on
a small sample" and "it works" are different claims for a pipeline whose failure mode is
cumulative rather than immediate.

\subsection{Limitations}
No single real dataset gives whole-organism 3D geometry for \textit{Arabidopsis}, just as
none exists for rice; every organ shape here is a citation-informed illustrative
approximation, not a reconstruction from measurement, and several structural parameters
(rosette leaf-blade curvature, inflorescence taper, flower whorl spacing) are not yet tied
to a specific measured value even where a general developmental-biology source is cited.
TraVA/Klepikova expression is linked to, not integrated into, the viewer, because its bulk
reuse terms could not be confirmed -- a more complete atlas would resolve that directly
with the maintainers rather than work around it indefinitely. OSD-38 and OSD-522, both
confirmed real and relevant, remain unprocessed. The spaceflight comparisons reported are
mean ratios from two studies with their own genotype and lighting confounds, not a
meta-analytic or statistically tested effect size, and should not be read as one. The
growth animation's timing rests on a single primary source \citep{Boyes2001} -- its
corroborating TAIR pages returned HTTP 403 to automated verification and were not
independently checked -- and its root-growth curve, having no Boyes-equivalent measurement
to follow, is a disclosed proxy (the rosette-growth fraction) rather than an independently
timed root phenotype.

\subsection{What would make this stronger}
Resolving TraVA's reuse terms directly, processing OSD-38 and OSD-522 to extend
organ-matched spaceflight coverage beyond root, and replacing illustrative structural
parameters with values measured from the cited primary sources (rather than only citing
the source's general conclusions) would each move this resource further from "a rice-atlas
pattern applied honestly" toward "a primary structural and functional atlas in its own
right."
